\section{\HMTLang{Reciprocidad diferencial y restitución de los canales constitutivos}{Differential reciprocity and recovery of constitutive channels}}
\label{delta22:iii:restitucion}
\HMTLang{Los canales del vacío, el funcional de Barbero--Immirzi y el registro dodecafásico pertenecen a una misma cadena de construcciones. El resultado siguiente reúne sus inversas: la velocidad normalizada conserva la componente común de la respuesta, mientras la lectura orientada de Barbero determina su reparto entre hojas. La composición recupera el par angular y, sobre la imagen generada con su sección dimensional, el coeficiente de acción. Conservando además las coordenadas regionales, alcanza la lectura periódica de \(K\).}{The vacuum channels, the Barbero--Immirzi functional, and the dodecaphasic record belong to one chain of constructions. The following result assembles their inverses: normalized speed retains the common component of the response, while the oriented Barbero reading determines its distribution between sheets. The composition recovers the angular pair and, on the generated image with its dimensional section, the action coefficient. Retaining the regional coordinates also gives the periodic reading of \(K\).}

\HMTLang{La construcción precedente de APP--TRIT--TPK y su estado enriquecido permanece como antecedente. Los parámetros empleados a continuación son las representaciones angulares ya obtenidas; las pruebas analíticas describen su realización y las operaciones de recuperación. Se mantienen explícitos los proyectores etiquetados \(P_\pm\), \(\mathcal R=P_+-P_-\), la orientación y el levantamiento angular.}{The preceding construction of APP--TRIT--TPK and its enriched state remains the antecedent. The parameters used below are the angular outputs already obtained; the analytic proofs describe their realization and recovery operations. The labelled projectors \(P_\pm\), \(\mathcal R=P_+-P_-\), orientation, and angular lift remain explicit.}

\subsection{\HMTLang{Potencial espectral y reciprocidad constitutiva}{Spectral potential and constitutive reciprocity}}
\HMTLang{En el cono \(x>|y|\), sean}{On the cone \(x>|y|\), let}
\[
 T=e^{-xI-y\mathcal R},\quad s_\pm=e^{-30(x\pm y)},\quad
 f(s)=\frac{1-s^3}{1-s^4},\quad
 V=(I-T^{90})(I-T^{120})^{-1}=r_+P_++r_-P_-,
\]
\[
 r_\pm=f(s_\pm),\quad
 \widehat c=(r_+r_-)^{-1},\quad \widehat Z=r_-/r_+,
 \quad c_\ell=\log\widehat c,\quad z_\ell=\log\widehat Z.
\]
\HMTLang{La invertibilidad de \(I-T^{120}\) se deduce de \(0<s_\pm<1\). El cociente de polinomios prolongado es \(f(s)=(1+s+s^2)/(1+s+s^2+s^3)\); sus extremos son \(1\) y \(3/4\).}{Invertibility of \(I-T^{120}\) follows from \(0<s_\pm<1\). The extended polynomial quotient is \(f(s)=(1+s+s^2)/(1+s+s^2+s^3)\); its endpoint values are \(1\) and \(3/4\).}
\begin{theorem}[\HMTLang{Potencial de la respuesta normalizada}{Potential of the normalized response}]
\HMTLang{La serie}{The series}
\begin{equation}
 \mathscr F(x,y)=\sum_{\sigma=\pm1}\sum_{n\ge1}
 \left[\frac{e^{-120n(x+\sigma y)}}{120n^2}
       -\frac{e^{-90n(x+\sigma y)}}{90n^2}\right]
 \label{delta22:iii:potencial}
\end{equation}
\HMTLang{define un potencial real analítico estrictamente cóncavo, con}{defines a strictly concave real-analytic potential, with}
\[
 \nabla\mathscr F=(c_\ell,z_\ell),\qquad
 \partial_y\log\widehat c=\partial_x\log\widehat Z.
\]
\HMTLang{La normalización es \(\mathscr F\to0\) cuando \(x-|y|\to\infty\).}{Its normalization is \(\mathscr F\to0\) as \(x-|y|\to\infty\).}
\end{theorem}
\begin{proof}
\HMTLang{En cada compacto del cono existe \(\epsilon>0\) con \(x\pm y\ge\epsilon\). Cada derivada de orden fijo de los sumandos queda dominada por una potencia de \(n\) multiplicada por \(e^{-90n\epsilon}\); la suma converge uniformemente. La derivación término a término y \(\sum_{n\ge1}z^n/n=-\log(1-z)\) dan las dos componentes indicadas. Para la concavidad, sean}{On each compact subset of the cone there is \(\epsilon>0\) with \(x\pm y\ge\epsilon\). Every fixed-order derivative of a summand is bounded by a power of \(n\) times \(e^{-90n\epsilon}\); the sum converges uniformly. Termwise differentiation and \(\sum_{n\ge1}z^n/n=-\log(1-z)\) give the stated components. For concavity, set}
\[
 g(w)=\log f(e^{-w}),\quad u=30(x+y),\quad v=30(x-y),
 \quad a=g'(u),\quad b=g'(v).
\]
\[
 g'(w)=\frac3{e^{3w}-1}-\frac4{e^{4w}-1}>0,
 \qquad
 D(c_\ell,z_\ell)=-30
 \begin{pmatrix}a+b&a-b\\a-b&a+b\end{pmatrix}.
\]
\HMTLang{La desigualdad se obtiene porque \(t/(e^{tw}-1)\) decrece estrictamente en \(t>0\): su derivada tiene numerador \(e^{tw}(1-tw)-1<0\). Los autovalores de la matriz son \(-60a\) y \(-60b\), ambos negativos, y su determinante es \(3600ab>0\). La simetría prueba la reciprocidad diferencial. La cota exponencial da el límite normalizador.}{The inequality follows because \(t/(e^{tw}-1)\) is strictly decreasing for \(t>0\): its derivative has numerator \(e^{tw}(1-tw)-1<0\). The matrix eigenvalues are \(-60a\) and \(-60b\), both negative, and its determinant is \(3600ab>0\). Symmetry proves differential reciprocity. The exponential bound gives the normalizing limit.}
\end{proof}
\HMTLang{El mismo momento espectral \(\mathcal D_2(z)=\sum_{n\ge1}z^n/n^2\) permite escribir \eqref{delta22:iii:potencial} mediante sus evaluaciones reales contractivas. El carácter orientado ya construido satisface \(\mathcal C_2(s)=\operatorname{Im}\mathcal D_2(is)\), \(\mathcal C_2(1)=G_{\rm Cat}\). La convergencia absoluta alcanza \(s=1\); para las series derivadas que dejan de converger allí se conserva el límite radial de Abel. Las evaluaciones reales y orientadas comparten el índice de profundidad y la ponderación, conservando sus argumentos respectivos. Bajo la amplificación \(T\otimes I_{9^m}\), la traza dividida por \(9^m\) conserva estas lecturas, pues cada autovalor se repite exactamente \(9^m\) veces. Esta identidad corresponde a esa amplificación declarada.}{The same spectral moment \(\mathcal D_2(z)=\sum_{n\ge1}z^n/n^2\) expresses \eqref{delta22:iii:potencial} through real contractive evaluations. The oriented character already constructed satisfies \(\mathcal C_2(s)=\operatorname{Im}\mathcal D_2(is)\), \(\mathcal C_2(1)=G_{\rm Cat}\). Absolute convergence extends to \(s=1\); for differentiated series that cease to converge there, the radial Abel limit is retained. Real and oriented evaluations share the depth index and weighting while retaining their respective arguments. Under amplification by \(T\otimes I_{9^m}\), the trace divided by \(9^m\) preserves these readings, since every eigenvalue is repeated exactly \(9^m\) times. This identity concerns the specified amplification.}

\begin{proposition}[\HMTLang{Sensibilidad orientada y convexidad}{Oriented sensitivity and convexity}]
\HMTLang{Para \(x>0\) fijo, \(y\mapsto c_\ell(x,y)\) es par y estrictamente convexa, con mínimo único en cero.}{For fixed \(x>0\), \(y\mapsto c_\ell(x,y)\) is even and strictly convex, with its unique minimum at zero.}
\end{proposition}
\begin{proof}
\[
 g''(w)=-\frac9{4\sinh^2(3w/2)}+\frac{16}{4\sinh^2(2w)}<0.
\]
\HMTLang{En efecto, \(t/\sinh(tw/2)\) decrece: su derivada tiene el signo de \(\sinh v-v\cosh v<0\). Como \(c_\ell=-g(u)-g(v)\), se obtiene \(\partial_y^2c_\ell=-900[g''(u)+g''(v)]>0\). La simetría bajo \(y\mapsto-y\) fija el mínimo. La expansión local es}{Indeed, \(t/\sinh(tw/2)\) decreases: its derivative has the sign of \(\sinh v-v\cosh v<0\). Since \(c_\ell=-g(u)-g(v)\), one obtains \(\partial_y^2c_\ell=-900[g''(u)+g''(v)]>0\). Symmetry under \(y\mapsto-y\) fixes the minimum. The local expansion is}
\[
 c_\ell=-2g(30x)-900g''(30x)y^2+O(y^4),\qquad
 z_\ell=-60g'(30x)y+O(y^3).
\]
\end{proof}
\HMTLang{La velocidad retiene la separación de canales a segundo orden; la impedancia conserva su orientación a primer orden. Evaluar el canal medio antes del lector elimina precisamente esas contribuciones.}{Speed retains channel separation at second order; impedance retains its orientation at first order. Evaluating the mean channel before applying the reader removes precisely those contributions.}

\subsection{\HMTLang{Inversión global mediante velocidad y Barbero--Immirzi}{Global inversion through speed and Barbero--Immirzi}}
\HMTLang{La inversa \(\psi=f^{-1}:(3/4,1)\to(0,1)\) está definida por la raíz única de}{The inverse \(\psi=f^{-1}:(3/4,1)\to(0,1)\) is defined by the unique root of}
\[
 rs^3+(r-1)(1+s+s^2)=0.
\]
\HMTLang{La unicidad se deduce de}{Uniqueness follows from}
\[
 f'(s)=-\frac{s^2(s^2+2s+3)}{(1+s+s^2+s^3)^2}<0.
\]
\HMTLang{Conservamos el funcional de retorno ya derivado de la incidencia hexada--octada:}{Retain the return functional already derived from hexad--octad incidence:}
\[
 \mathcal H(s)=\frac{12s^3}{1-s^9}-\frac{s^4}{1-s^{12}}
              -\frac{(\log s)^2}{20\pi},\qquad F_{\rm B}=\mathcal H\circ\psi.
\]
\begin{theorem}[\HMTLang{Restitución constitutiva orientada}{Oriented constitutive recovery}]
\HMTLang{La aplicación}{The map}
\[
 (r_+,r_-)\longmapsto
 \left(\widehat c=\frac1{r_+r_-},\quad
 \gamma=F_{\rm B}(r_-)-F_{\rm B}(r_+)\right)
\]
\HMTLang{es un difeomorfismo real analítico de \((3/4,1)^2\) sobre \((1,16/9)\times\mathbb R\). Esta afirmación describe el dominio analítico del lector; la imagen efectivamente generada conserva, además, las compatibilidades del TPK.}{is a real-analytic diffeomorphism from \((3/4,1)^2\) onto \((1,16/9)\times\mathbb R\). This statement describes the analytic domain of the reader; the actually generated image additionally retains the TPK compatibility conditions.}
\end{theorem}
\begin{proof}
\[
 \mathcal H'(s)=\frac{36s^2(1+2s^9)}{(1-s^9)^2}
 -\frac{4s^3(1+2s^{12})}{(1-s^{12})^2}
 -\frac{\log s}{10\pi s}>0.
\]
\HMTLang{El cociente del primer término positivo por la magnitud del segundo es}{The ratio of the first positive term to the magnitude of the second is}
\[
 \frac9s\frac{1+2s^9}{1+2s^{12}}
 \left(\frac{1-s^{12}}{1-s^9}\right)^2>9,
\]
\HMTLang{y el último sumando también es positivo. Los extremos son \(\mathcal H(s)\sim-(\log s)^2/(20\pi)\) en cero y \(\mathcal H(s)\sim5/[4(1-s)]\) en uno. Por tanto \(F_{\rm B}'<0\), con límites \(+\infty\) en \(3/4\) y \(-\infty\) en \(1\).}{and the last summand is also positive. The endpoint asymptotics are \(\mathcal H(s)\sim-(\log s)^2/(20\pi)\) at zero and \(\mathcal H(s)\sim5/[4(1-s)]\) at one. Hence \(F_{\rm B}'<0\), with limits \(+\infty\) at \(3/4\) and \(-\infty\) at \(1\).}

\HMTLang{Para \(c=\widehat c\) fijo, sea \(z=\log\widehat Z\). Entonces}{For fixed \(c=\widehat c\), set \(z=\log\widehat Z\). Then}
\[
 r_+=c^{-1/2}e^{-z/2},\quad r_-=c^{-1/2}e^{z/2},\quad
 |z|<a(c)=\min\{\log c,\log(16/(9c))\}.
\]
\HMTLang{La función \(\mathcal G_c(z)=F_{\rm B}(r_-)-F_{\rm B}(r_+)\) satisface}{The function \(\mathcal G_c(z)=F_{\rm B}(r_-)-F_{\rm B}(r_+)\) satisfies}
\[
 \mathcal G_c'(z)=\tfrac12
 [r_-F_{\rm B}'(r_-)+r_+F_{\rm B}'(r_+)]<0.
\]
\HMTLang{En el extremo superior, \(r_-\to1\), o \(r_+\to3/4\), o ambos, y \(\mathcal G_c\to-\infty\). En el inferior tiende a \(+\infty\). El teorema del valor intermedio y la monotonía dan una única solución para cada \(\gamma\). La derivada no nula proporciona la inversa analítica por el teorema de la función implícita. La descripción por el mínimo de \(a(c)\) sólo delimita la frontera; el argumento se realiza en su dominio abierto.}{At the upper endpoint, \(r_-\to1\), or \(r_+\to3/4\), or both, and \(\mathcal G_c\to-\infty\). At the lower endpoint it tends to \(+\infty\). The intermediate value theorem and monotonicity give a unique solution for each \(\gamma\). The nonzero derivative gives the analytic inverse by the implicit function theorem. The minimum in \(a(c)\) only describes the boundary; the argument takes place in its open domain.}
\end{proof}

\subsection{\HMTLang{Recuperación angular, acción y registro periódico}{Angular recovery, action, and periodic record}}
\HMTLang{Una vez obtenidos \(r_\pm\), la inversión cúbica da \(s_\pm\). Con el levantamiento y la orientación conservados,}{Once \(r_\pm\) have been obtained, cubic inversion gives \(s_\pm\). With the lift and orientation retained,}
\[
 x=-\frac1{60}\log(s_+s_-),\quad y=\frac1{60}\log\frac{s_-}{s_+},
 \quad A_{\deg}=-\frac3\pi\log(s_+s_-),\quad
 C^*_{\deg}=\frac3\pi\log\frac{s_-}{s_+}.
\]
\HMTLang{Sobre la sección positiva generada \(A_{\deg}=1000\alpha\), \(C^*_{\deg}=2(\eta_{\rm ret}+\alpha)\), se recupera}{On the generated positive section \(A_{\deg}=1000\alpha\), \(C^*_{\deg}=2(\eta_{\rm ret}+\alpha)\), one recovers}
\[
 \alpha=A_{\deg}/1000,\qquad
 \eta_{\rm ret}=C^*_{\deg}/2-\alpha,\qquad
 \hbar_{\rm ret}=s_0\eta_{\rm ret},\qquad s_0:=10^{-34}\mathcal U_S.
\]
\HMTLang{La escala \(s_0\) y los proyectores son antecedentes conservados. En la orientación conjugada se transporta también la etiqueta de acción. La recuperación es una inversa sobre la imagen correspondiente.}{The scale \(s_0\) and the projectors are retained antecedents. Under conjugate orientation, the action label is transported as well. Recovery is an inverse on the corresponding image.}

\begin{proposition}[\HMTLang{Composición con la lectura periódica de doce bloques}{Composition with the twelve-block periodic reading}]
\HMTLang{La lectura periódica prolonga los doce bloques de \(K\), mientras las coordenadas regionales conservan su generación coinductiva completa. Con \(\alpha_A\) como coordenada de la clausura aritmética previamente definida,}{The periodic reading prolongs the twelve blocks of \(K\), while the regional coordinates retain their complete coinductive generation. With \(\alpha_A\) denoting the previously defined arithmetic-closure coordinate,}
\[
 \kappa_{\rm per}=\pi_{\rm HMT}+e_{\rm HMT}-\varphi_{\rm HMT}-4-\alpha_A,
 \quad B=1000,\quad N_K=(B^{12}-1)\kappa_{\rm per}.
\]
\HMTLang{Sobre la imagen del registro generado, \(N_K\) es entero y sus doce bloques, con origen fijo y ceros iniciales conservados, se recuperan exactamente por}{On the image of the generated record, \(N_K\) is an integer and its twelve blocks, with fixed origin and retained leading zeros, are recovered exactly by}
\[
 K_j=\left\lfloor\frac{N_K}{B^{12-j}}\right\rfloor
 -B\left\lfloor\frac{N_K}{B^{13-j}}\right\rfloor,
 \qquad 1\le j\le12.
\]
\end{proposition}
\begin{proof}
\HMTLang{La división euclídea sucesiva de \(N_K\) en base \(B\) determina de manera única \(K_j\in\{0,\ldots,B-1\}\) y produce \(N_K=\sum_{j=1}^{12}K_jB^{12-j}\). La repetición periódica de esos bloques suma la serie geométrica \(N_K\sum_{m\ge1}B^{-12m}=N_K/(B^{12}-1)\). Las dos composiciones son inversas. La lectura finita \(N_K/B^{12}\) conserva otra normalización, distinguida de la periódica. Con datos aproximados, la recuperación exige un intervalo certificado para \(N_K\) que contenga un único entero.}{Successive Euclidean division of \(N_K\) in base \(B\) uniquely determines \(K_j\in\{0,\ldots,B-1\}\) and gives \(N_K=\sum_{j=1}^{12}K_jB^{12-j}\). Periodic repetition of these blocks sums the geometric series \(N_K\sum_{m\ge1}B^{-12m}=N_K/(B^{12}-1)\). Both compositions are inverse. The finite reading \(N_K/B^{12}\) retains a different normalization from the periodic one. With approximate data, recovery requires a certified interval for \(N_K\) containing a unique integer.}
\end{proof}
\HMTLang{Esta composición reúne la restitución constitutiva con la recuperación del registro ya construido. La selección excepcional y temporal determina \(K\) previamente; la inversa explica qué información permiten restituir sus lectores. El alcance es el registro ordenado de doce coordenadas con los antecedentes indicados.}{This composition assembles constitutive restitution with recovery of the already constructed record. Exceptional and temporal selection determines \(K\) beforehand; the inverse explains which information its readers can restore. Its scope is the ordered twelve-coordinate record with the stated antecedents.}

\subsection{\HMTLang{Cotas de estabilidad en el dominio contractivo}{Stability bounds in the contractive domain}}
\HMTLang{Para \(w>0\), escribamos \(a(w)=g'(w)\). Con \(s=e^{-w}\),}{For \(w>0\), write \(a(w)=g'(w)\). With \(s=e^{-w}\),}
\[
 a(w)=\frac{s^3(s^2+2s+3)}{(1+s+s^2)(1+s+s^2+s^3)},
 \qquad \tfrac12e^{-3w}<a(w)<3e^{-3w}.
\]
\HMTLang{Las cotas se obtienen al comparar los polinomios positivos del cociente; \(g''<0\) demuestra su decrecimiento. En un dominio convexo donde \(0<m_0\le30(x\pm y)\le M_0<\infty\), la aplicación \(\mathscr C=(c_\ell,z_\ell)\) satisface, con la norma euclídea,}{The bounds follow by comparing the positive polynomials in the quotient; \(g''<0\) proves its decrease. On a convex domain where \(0<m_0\le30(x\pm y)\le M_0<\infty\), the map \(\mathscr C=(c_\ell,z_\ell)\) satisfies, in the Euclidean norm,}
\[
 60a(M_0)\|p-q\|\le\|\mathscr C(p)-\mathscr C(q)\|
 \le60a(m_0)\|p-q\|.
\]
\HMTLang{Para la cota superior se integra el jacobiano sobre el segmento. Para la inferior se aplica el producto escalar con \(p-q\) al jacobiano negativo definido y luego Cauchy--Schwarz. La norma inducida de la inversa diferencial es \(1/[60\min\{a(u),a(v)\}]\). En contracción profunda puede crecer exponencialmente, incluso cuando el número de condición relativo permanece próximo a uno. La unicidad exacta y la robustez frente al error quedan, por tanto, especificadas por criterios diferentes y compatibles.}{For the upper bound, integrate the Jacobian along the segment. For the lower bound, take the scalar product with \(p-q\), use negative definiteness of the Jacobian, and then apply Cauchy--Schwarz. The induced norm of the inverse differential is \(1/[60\min\{a(u),a(v)\}]\). In deep contraction it can grow exponentially even when the relative condition number remains close to one. Exact uniqueness and robustness to error are thus specified by distinct, compatible criteria.}
